\documentclass[11pt]{article}
\usepackage[margin=1in]{geometry}
\usepackage{amsmath,amssymb,booktabs,longtable}
\usepackage[hidelinks]{hyperref}
\newcommand{\dd}{\mathrm d}
\newcommand{\lp}{\lambda_{HX}}
\newcommand{\lt}{\lambda_{\rm tot}}
\title{The particle bridge beyond zero-temperature stability\\
Thermal phases, charge-resolved scattering, and the next kinetic equations}
\author{Zimmerman programme: CPB-TC-2026-10-16}
\date{16 October 2026}
\begin{document}
\maketitle
\begin{abstract}
This continues the explicit scalar Higgs/carrier construction at checkpoint
CPB-2026-10-16. A leading high-temperature Higgs-bath comparison finds an
intermediate instability of the source field at the abundance-calibrated
1-MeV benchmark. An exact sufficient inequality avoids dark condensation in
the stated class of positive bath-mass comparisons; a separately dated
mass-family continuation satisfies it. Vacuum tree amplitudes show that the
portal-produced neutral mixture scatters more strongly than the previous
one-channel estimate. Time dilation also invalidates a prompt-source-decay
shortcut in a separate vacuum control. Neither a physical thermal history,
a relic abundance, the galaxy law, nor full theory closure is established.
\end{abstract}

\section{Continuation contract and model}
The previous goal turn was progress: it produced an explicit reciprocal
scalar potential and identified unresolved thermal/kinetic conditions.
The actual checkout remains
\texttt{370dfb58006e330f210f5e306297e822188e37ba}, with shared uncommitted
work. This continuation reads the actual prior report, its current run 003,
and the October 15 work order. No frozen theory, seal, prior empirical score,
or other lane is modified. The new authoritative computation is
\texttt{thermal\_charge\_run\_002/manifest.json}. This is self-review.

Use $\hbar=c=k_B=1$, $(+---)$, and the same complex scalars $\Phi,X$ with
charges $1,2$. Write $x=|X|$, $n=|\Phi|^2$, $t=H^\dagger H$, and
$w_H=t-v_H^2/2$. The previously specified potential is
\begin{align}
 V_J={}&m^2|\Phi|^2+M^2|X|^2
 +g_c(X^*\Phi^2+X\Phi^{*2})+\lt|\Phi|^4\\
 &+\lambda_Hw_H^2+\lp w_H|X|^2+\lambda_X|X|^4,\nonumber\\
 M^2={}&4m^2(1+u_*^2),\quad \Delta_c=4m^2u_*^2,\quad
 \lt=\frac{g_c^2}{\Delta_c}+\lambda_{\rm rep}.
\end{align}
The positive-square carrier identity and
$\lambda_X\ge\lp^2/(4\lambda_H)$ establish the previous global
zero-temperature bound. They do not prove that the same origin is stable
at every temperature. In particular,
\begin{equation}
 M_0^2=M^2-\lp v_H^2/2
\end{equation}
is the source curvature before the Higgs obtains its vacuum expectation value.

The benchmark fixes $m=10^{-3}$ GeV, $u_*=.02$, $r_\lambda=2$,
$\sigma_0/m=.3\ \mathrm{cm^2/g}$, $\lambda_X=10^{-18}$ and
$\lp=2.21149\,10^{-10}$. The portal was calibrated in the earlier
simplified production control to $\Omega_c h^2=.12$; it is not a derived
abundance. $g_c\simeq2.168\,10^{-6}$ GeV and $\lt\simeq.00881$ follow
from these stipulated microscopic inputs.

\section{A thermal comparison that does not assume a dark heat bath}
Assume the Standard Model bath is thermal, initially no populated dark bath,
and retain only its leading high-temperature real scalar self-energies.
Specify the comparison free energy
\begin{equation}
 V_T=V_J+c_H T^2t+\frac{\lp T^2}{6}x^2,
 \qquad
 c_H=\frac{3g_2^2+g_Y^2}{16}+\frac{\lambda_H}{2}+\frac{y_t^2}{4}.
\label{eq:VT}
\end{equation}
We do not insert $\lt T^2$ or $\lambda_XT^2$ as dark thermal masses:
the dark distributions have not been shown to equilibrate at the visible
temperature. Those corrections would depend on the actual populated states.

The portal tadpole normalization is independently fixed by
\begin{equation}
 \langle H^\dagger H\rangle_T=
 \frac12\sum_{a=1}^4\langle h_a^2\rangle_T=\frac{T^2}{6}
\end{equation}
for four massless real Higgs components with
$\langle h_a^2\rangle_T=T^2/12$. This is the same normalization as
Seller, Sz\'ep and Tr\'ocs\'anyi, Appendix C, equations (C.11)--(C.12)
\cite{seller}: their complex singlet is our $X$, their doublet is our $H$,
and their quartic mixing coefficient is our $\lp$. Their reverse singlet-bath
contribution to the Higgs mass is omitted here because there is no assumed
dark thermal population. The visible-bath diagram remains present even
when the dark field is out of equilibrium.

With illustrative fixed $g_2=.65$, $g_Y=.36$, $y_t=.94$,
$\lambda_H=m_h^2/(2v_H^2)$, $m_h=125$ GeV and $v_H=246$ GeV,
$c_H=.372768$. The comparison Higgs origin changes stability at
\begin{equation}
 T_H=\sqrt{\lambda_Hv_H^2/c_H}=144.77\ \mathrm{GeV}.
\end{equation}
This is a mean-field control scale, not a new prediction for the physical
Standard Model crossover. Running, thermal resummation, field-dependent
thermal integrals, imaginary parts in unstable regions, and non-equilibrium
dark backreaction have not been solved. High-temperature formulas must not
be extrapolated to a few GeV as physical decay-threshold predictions.
Phase and thermal-rate effects in freeze-in are known issues \cite{deromeri}.

\section{Analytic elimination and the baseline instability}
Minimize over the relative scalar phase so that the cubic term is
$-2g_cxn$. At each fixed $x\ge0$ the two other invariants minimize exactly at
\begin{equation}
 n_*(x)=\max\!\left(\frac{2g_cx-m^2}{2\lt},0\right),\quad
 t_*(x,T)=\max\!\left(\frac{\lambda_Hv_H^2-c_HT^2-\lp x^2}
 {2\lambda_H},0\right).
\label{eq:elim}
\end{equation}
Thus the entire minimization of the stipulated homogeneous free energy is
one-dimensional and piecewise quartic. The boundaries are
$x=m^2/(2g_c)$ and, when real,
$x=\sqrt{(\lambda_Hv_H^2-c_HT^2)/\lp}$.
Interior stationary points obey
\begin{equation}
 4a x^3+2b x+c=0.
\end{equation}
For $t_*>0$,
$a=\lambda_X-\lp^2/(4\lambda_H)$ and
$b=M^2+\lp[1/6-c_H/(2\lambda_H)]T^2$.
For $t_*=0$, $a=\lambda_X$ and $b=M_0^2+\lp T^2/6$.
If $n_*>0$, replace $b$ by $b-g_c^2/\lt$ and set
$c=g_cm^2/\lt$; otherwise $c=0$.
Checking all admissible roots, branch boundaries and $x=0$ suffices since
the quartic coefficients are positive. The implementation enumerates those
roots and separately checks a logarithmic amplitude grid.

In the unbroken Higgs branch the source's linear curvature is
\begin{equation}
 M_T^2=M_0^2+\lp T^2/6.
\end{equation}
At the benchmark $M_0^2=-2.68992\,10^{-6}$ GeV$^2$. At $T=200$ GeV,
\begin{equation}
 M_T^2=-1.21559\,10^{-6}\ \mathrm{GeV}^2<0.
\end{equation}
Both visible Higgs and daughter have positive quadratic curvatures at that
origin, so this source instability is an actual negative scalar direction
of the specified comparison, not an artifact of minimizing them out.
A canonical source fluctuation satisfies, locally,
\begin{equation}
 \ddot{\delta X}+3H\dot{\delta X}
 +(k^2/a^2+M_T^2)\delta X=0.
\end{equation}
For long wavelengths $\sqrt{-M_T^2}/H=1.9621\,10^{10}$ with the stated
radiation-era $H$. The exact classical field at precisely zero remains at
zero without a seed; the result is linear instability to admitted nonzero
fluctuations, not spontaneous production proved by a deterministic zero solution.

The stipulated high-temperature source-origin instability occupies
\begin{equation}
 T_H<T<\sqrt{-6M_0^2/\lp}=270.15\ \mathrm{GeV}.
\end{equation}
The exact lower edge of the full comparison's dark phase can extend below
$T_H$; it is not asserted here as a physical transition temperature.
Varying $c_H$ over $.3,.4,.5$ and the source tadpole coefficient over
$\lp/12,\lp/6,\lp/3$ preserves a nonempty unbroken-Higgs instability
interval in all nine stipulated cases. This is coefficient sensitivity,
not a precision resummed calculation.

At 200 GeV the minimized comparison gives
\begin{equation}
 |X|=7.79785\,10^5\ \mathrm{GeV},\quad
 |\Phi|=13.8502\ \mathrm{GeV},\quad H=0,\quad
 \Delta V_T=-3.69742\,10^5\ \mathrm{GeV}^4.
\end{equation}
These are homogeneous mean-field extrema and a free-energy difference,
not a simulated condensate, a released-energy budget, or a relic yield.
Both dark amplitudes are nonzero; the vacuum $U(1)$ is broken on that
comparison branch and the original uncondensed particle spectrum cannot
be reused there. Whether the universe follows this branch, fragments,
restores the symmetry or produces acceptable relics is unresolved.
The earlier zero-temperature positivity proof remains correct. What fails
is promoting it to an unbroken dilute thermal production history.

\section{A sufficient condition for avoiding the dark phase}
There is a useful analytic continuation beyond that failed shortcut.
Suppose $\lp,\lambda_X,\lt>0$, the only dark bath correction is a
nonnegative source-mass correction, and
\begin{equation}
 \boxed{M_0^2\ge g_c^2/\lt.}
\label{eq:sufficient}
\end{equation}
For all $x,n,t\ge0$ the dark part before that bath correction is
\begin{align}
 &m^2n+M_0^2x^2-2g_cxn+\lt n^2+\lp tx^2+\lambda_Xx^4\\
 &\quad=\lt\left(n-\frac{g_cx}{\lt}\right)^2
 +m^2n+\left(M_0^2-\frac{g_c^2}{\lt}\right)x^2
 +\lp tx^2+\lambda_Xx^4\ \ge0.\nonumber
\end{align}
The relative phase can only raise this minimum. For each Higgs field
configuration, $X=\Phi=0$ attains the dark minimum. Therefore no lower
dark-condensed equilibrium exists in this whole specified class, even if
the visible thermal potential is replaced by an arbitrary bounded function
of $t$ with an attained minimum. This is a sufficient inequality, not a
necessary condition, and not a theorem including uncomputed dark thermal
or quantum corrections.

To execute a changed continuation, hold $u_*,r_\lambda,\sigma_0/m$ and the
earlier illustrative final abundance fixed while varying $m$. The prescribed
relations are
\begin{equation}
 \lambda_{\rm rep}=\sqrt{8\pi m^3(\sigma_0/m)},\quad
 g_c^2=\lambda_{\rm rep}\Delta_c/r_\lambda,\quad
 \lp(m)\propto m^{-1/2}(1-4M^2/m_h^2)^{-1/4}.
\end{equation}
The last is abundance calibration within the old decay-yield approximation;
it is not a cosmologically verified prescription. Substitution into
equation~\eqref{eq:sufficient} gives the illustrative sufficient floor
$m\simeq1.22840$ MeV. This number is neither a derived mass nor an observed
particle bound.

\begin{center}
\begin{tabular}{r r l}
\toprule $m$ (MeV) & $M_0^2$ (GeV$^2$) & Condition~\eqref{eq:sufficient}\\\midrule
0.75 & $-5.47580\,10^{-6}$ & not met\\
1.00 & $-2.68992\,10^{-6}$ & not met\\
1.25 & $2.67422\,10^{-7}$ & met\\
2.00 & $1.12748\,10^{-5}$ & met\\
10.0 & $3.98044\,10^{-4}$ & met\\\bottomrule
\end{tabular}
\end{center}
The 1.25-, 2- and 10-MeV rows have no lower dark-condensed state at the
nine tested temperatures; the all-temperature sufficient statement follows
from the square identity, not that finite grid. These are separately
specified mass-family candidates, not a rescue declared for the old point.
They still require thermal rates, full kinetics, empirical tests and the
connection to gravity. More massive candidates can also change independent
laboratory/cosmological constraints, which have not been scored here.

\section{The neutral population has two distinct scattering channels}
The Higgs portal produces $XX^*$ and hence equal positive/negative daughter
populations in the stated charge-symmetric scenario. The conserved charge
is a difference, not their total rest mass. The vacuum tree vertices from
$-\lt|\Phi|^4-g_c(X^*\Phi^2+\mathrm{h.c.})$ are $-4i\lt$ and $-2ig_c$.
With on-shell external daughters and invariant $s$,
\begin{align}
 \mathcal M_{++}(s)&=-4\left[\lt+\frac{g_c^2}{s-M^2}\right],\\
 \mathcal M_{+-}(u)&=-4\left[\lt+\frac{g_c^2}{u-M^2}\right].
\end{align}
Same-charge scattering has an $s$-channel source; the crossed
particle--antiparticle channel has $u$-channel exchange. There is no
same-vertex $X$ $t$-channel because it would violate charge assignment.
These are tree amplitudes; the physical near-pole same-charge channel needs
a width and inclusive accounting. The crossed channel has no such pole for
the physical elastic $u\le0$ used here.

For COM momentum $p$, $s=4(m^2+p^2)$ and
$u=-2p^2(1+\cos\theta)$. For distinguishable opposite-charge daughters,
\begin{equation}
 \sigma_{+-}(s)=\frac{1}{2\pi s}\int_{-1}^1
 \left[\lt-\frac{g_c^2}{M^2+2p^2(1+z)}\right]^2\dd z.
\end{equation}
For identical same-charge final particles the $1/2!$ factor instead gives
\begin{equation}
 \sigma_{++}(s)=\frac{1}{2\pi s}
 \left[\lt+\frac{g_c^2}{s-M^2}\right]^2.
\end{equation}
At threshold,
\begin{align}
 \sigma_{++}(0)&=\frac{\lambda_{\rm rep}^2}{8\pi m^2}\equiv\sigma_0,\\
 \sigma_{+-}(0)&=\frac{(\lt-g_c^2/M^2)^2}{4\pi m^2},\\
 \frac{\sigma_{+-}(0)}{\sigma_0}
 &=2\left[1+\frac{1}{r_\lambda(1+u_*^2)}\right]^2=4.49880056.
\end{align}
The threshold rate per test daughter in an equal mixture of total density
$n_{\rm tot}$ and equal isotropic speed distributions is
\begin{equation}
 \Gamma_{\rm el}=\frac{n_{\rm tot}}{2}
 (\sigma_{++}+\sigma_{+-})\langle v_{\rm rel}\rangle.
\end{equation}
It is $2.74940028$ times the old screen $n_{\rm tot}\sigma_0
\langle v_{\rm rel}\rangle$. Therefore the old 1-keV monochromatic
collisionless control has $\Gamma_{\rm el}/H\simeq3776.90$, rather than
1373.72, in the threshold limit. This refines a hypothetical control and
does not supply the evolved spectrum. These are total collision cross
sections, not automatically a viscosity or momentum-transfer cross section
to insert into an empirical halo bound. Angular weighting needs its own
definition, especially for identical particles.

Nine exact rational threshold identities and numerical angular integrations
at $p/m=0,.00025,.003,.01,.05$ check normalization and continuity. They do
not compute the dense-medium inelastic collision operator.

\section{A short resting lifetime does not prove prompt cosmological decay}
The earlier isolated source lifetime is about $3.52\,10^{-13}$ seconds.
For a monochromatic Higgs initially at rest, $\gamma_{X,p}=m_h/(2M)$,
and the production-era decay rate is reduced by that boost. Keeping vacuum
masses and constant radiation degrees of freedom solely as a separate control,
\begin{equation}
 \frac{\Gamma_X}{\gamma_{X,p}H_p}=.0245331\quad(T_p=m_h/3).
\end{equation}
This is not prompt on the production Hubble scale. If $r=a/a_p$,
$u_{X,p}=\sqrt{\gamma_{X,p}^2-1}$ and $H=H_p/r^2$, the survival law is
\begin{equation}
 P_X(r)=\exp[-\tau(r)],\qquad
 \tau(r)=\frac{\Gamma_X}{H_p}\int_1^r
 \frac{s^2}{\sqrt{s^2+u_{X,p}^2}}\dd s.
\end{equation}
Numerical quadrature gives $\tau=1$ at $r=4.97701$ or $T=8.37182$ GeV,
with the source still highly relativistic. These numbers omit thermal
mass shifts and dark collisions and are not a physical decay epoch for
the unstable thermal benchmark.
While the source is ultrarelativistic, delayed decay approximately preserves
the mean comoving daughter momentum, since $aE_X$ is approximately constant.
Thus this delay alone does not overturn the earlier kinematic redshift
estimate. It does invalidate using a prompt cascade in the rate evolution.

\section{An explicit number-changing amplitude, not just an allowed charge}
The stronger check is the actual reaction
\begin{equation}
 2\Phi+2\Phi^*\longrightarrow\Phi+\Phi^*.
\end{equation}
Take all four incoming particles at rest. Each final particle has energy
$2m$ and momentum magnitude $\sqrt3m$. In the specified scalar theory,
with $\lambda_{H\Phi}=0$ at tree level, the complete six-daughter tree
graphs have three topologies: nine graphs with two daughter quartics,
18 with a daughter quartic and two cubic source vertices, and nine with
four cubic source vertices. The source quartic and Higgs portal cannot
occur in a six-daughter tree at this order: with trivalent/quartic vertices
$n_3+2n_4=4$, those choices leave too few external daughters or unmatched
Higgs/source legs. This does not omit them from loops or other processes.

The sum can be reconstructed without a current recursion. In units $m=1$,
write $d=\Delta_c/m^2$, $L=M^2/m^2+2=6+d$ and $\widehat g=g_c/m$.
An external same-charge pair has source denominator $-d$ if both are
incoming or $-L$ if one is a crossed outgoing particle. Daughter triple
denominators are $8$ or $-4$. For the three topologies the sums of inverse
propagator products are respectively
\begin{equation}
 S_0=-\frac34,\qquad S_2=\frac{3}{2L},\qquad
 S_4=-\frac{1}{4d^2}+\frac{1}{2dL}-\frac{1}{L^2}.
\end{equation}
For $S_0$, four triples give denominator $8$ and five give $-4$.
For $S_2$, endpoint pairs of invariant $4$ cancel after summing their
three allowed attachments; the other pairs give the displayed $3/(2L)$.
For $S_4$, the two endpoint pairs can be incoming/incoming once,
incoming/crossed in four cases, or crossed/crossed in four cases,
giving the three displayed terms. The products of vertex factors are
$16\lt^2$, $16\lt\widehat g^2$, and $16\widehat g^4$, with the
common amplitude sign $-1$. Restoring units,
\begin{equation}
 \mathcal M_{4\to2}=\frac1{m^2}\left[
 12\lt^2-\frac{24\lt\widehat g^2}{L}
 +\widehat g^4\left(\frac4{d^2}-\frac8{dL}+\frac{16}{L^2}\right)
 \right].
\end{equation}
Using $\widehat g^2=\lambda_{\rm rep}d/r_\lambda$ and
$z=d/L\in(0,1)$ gives a useful universal simplification:
\begin{equation}
 \boxed{\mathcal M_{4\to2}
 =\frac{\lambda_{\rm rep}^2}{m^2}
 \left[12+\frac{24(1-z)}{r_\lambda}
 +\frac{16(1-z)^2}{r_\lambda^2}\right].}
\label{eq:six}
\end{equation}
Consequently this vacuum tree amplitude is strictly nonzero for
$m,\Delta_c,g_c^2,\lambda_{\rm rep}>0$. The bracket is greater than 12.
This is a statement about the specified tree theory throughout that
parameter family, not a statement that chemical equilibration occurs.
It also does not establish the in-medium amplitude near a changing phase.

At $u_*^2=1/2500$, $r_\lambda=2$, its bracket is exactly
$393885012/14070001\simeq27.9947$. The computation uses exact rational
momenta and polynomials in $g_c/m$. A labeled off-shell-current recursion
and a separate direct enumeration of the 36 graph topologies agree at
12 rational parameter points. Re-rooting the external labels, the analytic
four-point amplitudes, and the pure-quartic six-point limit supply controls.
The positive closed form proves nonvanishing; the finite checks only
verify those implementations. This remains self-review.

For a dilute nonrelativistic uncorrelated gas, the zero-incoming-momentum
event rate of this one channel is
\begin{equation}
 R_{2+2-\to+-}=\frac{\sqrt3\,|\mathcal M_{4\to2}|^2}
 {1024\pi m^4}n_+^2n_-^2.
\end{equation}
The factors are the initial $1/(2!2!)$, four initial measures $1/(2m)$,
and the distinguishable final two-body phase space $\sqrt3/(16\pi)$.
Since two particles are lost per event, a symmetric gas of total density
$n$ has one-channel loss coefficient
\begin{equation}
 \Gamma_N\equiv\frac{2R}{n}
 =\frac{\sqrt3\,|\mathcal M_{4\to2}|^2 n^3}{8192\pi m^4}.
\end{equation}
Using the earlier stipulated density in a threshold control gives
$\Gamma_N/H=1.07\,10^{-28}$ at visible $T=1$ keV and
$1.24\,10^{-6}$ at 1 MeV. The latter has control $p/m\simeq.35$ and
cannot be read as a precision threshold rate, especially with the narrow
source scale $u_*=.02$. Neither number checks relativistic early production,
inverse gains, other charge channels or medium effects. They show why
``number change exists'' and ``number change is fast'' must be kept separate.

\section{The kinetic system now required}
The actual species are $\Phi,\Phi^*,X,X^*$, the visible bath, and any
coherent phase/response fields. In an uncondensed quasiparticle phase,
\begin{equation}
 (\partial_t-Hp\partial_p)f_i=C_i^{\rm portal}
 +C_i^{1\leftrightarrow2}+C_i^{2\leftrightarrow2}
 +C_i^{\rm number\ change}.
\end{equation}
An isotropic reduction still retains momentum dependence. With
$\dd\Pi_j=\dd^3p_j/[(2\pi)^32E_j]$, a bosonic elastic term has the form
\begin{align}
 C_a(p_1)={}&\frac{1}{2E_1 S_f}\int\dd\Pi_2\dd\Pi_3\dd\Pi_4
 (2\pi)^4\delta^4(p_1+p_2-p_3-p_4)|\mathcal M|^2\\
 &\ \times[f_3f_4(1+f_1)(1+f_2)-f_1f_2(1+f_3)(1+f_4)],\nonumber
\end{align}
where for the displayed tagged-daughter elastic channels $S_f=2!$ for
same-charge identical final daughters and $S_f=1$ for opposite-charge
final daughters. Other processes need their own tagged-species multiplicities.
The source collision term is
\begin{align}
 C_X(p_X)={}&\frac{1}{2E_X\,2!}\int\dd\Pi_1\dd\Pi_2
 (2\pi)^4\delta^4(p_X-p_1-p_2)\,4g_c^2\\
 &\ \times[f_1f_2(1+f_X)-f_X(1+f_1)(1+f_2)].\nonumber
\end{align}
The tagged daughter gain/loss term has the corresponding factor of two
relative to this event count. $X^*$ and the negative-charge daughters have
the conjugate terms. This is how inverse conversion, stimulation and charge
balance enter; a fixed decay rate without inverse terms is insufficient.

The classical global symmetry conserves
\begin{equation}
 q_\Phi=1,\ q_{\Phi^*}=-1,\ q_X=2,\ q_{X^*}=-2,\qquad
 \sum_i q_i\int\frac{\dd^3p}{(2\pi)^3} C_i=0.
\end{equation}
It does not conserve the total number of positive plus negative daughters.
The preceding section proves a nonzero mixed-charge six-point tree amplitude.
Its full relativistic/in-medium collision kernel and the other number-changing
channels have not been calculated, so neither total negligibility nor
chemical equilibrium may be asserted. In chemical equilibrium this reaction
imposes $\mu_\Phi+\mu_{\Phi^*}=0$; a charge-symmetric state then has zero
chemical potential for both. Elastic equilibration alone does not impose that.
The earlier all-positive charge-energy protection cannot silently be used
as a conservation theorem for a neutral population's rest abundance.
In particular, the earlier regulated pressure bound was formulated for the
positive-charge nonrelativistic sector with $d=\langle D_c\rangle/(m\langle Q\rangle)$.
For a charge-symmetric cosmological population $Q=0$ this quantity is
undefined and the bound is not a proof of continued coldness. A new
particle/antiparticle energy and pressure calculation is required. The
reaction above conserves total energy by converting some rest mass into
final kinetic energy; it supplies no energy from nothing.

For fixed masses, the sum of collision energy moments vanishes when every
participating bath/species is included. If quasiparticle masses depend on
the thermal/condensate background, each density also has a mass-work term:
\begin{equation}
 \dot\epsilon_i+3H(\epsilon_i+p_i)
 =\int\frac{\dd^3p}{(2\pi)^3}E_i C_i
 +\int\frac{\dd^3p}{(2\pi)^3}\dot E_i f_i.
\end{equation}
Mean-field and bath stresses must supply the reciprocal work without
double-counting the self-energy. Minima of $V_T$ are not by themselves
conserved mechanical energy or a production calculation.

This also sharpens the dark-energy question. A thermal scalar free-energy
offset is not automatically the observed cosmic vacuum. Any proposal that
net energy enters or leaves that vacuum must derive its $Q^\mu$ and obey
the previous checkpoint's exchange--acceleration consistency relation.
No equation here fixes the vacuum magnitude, $\kappa$, or the polarization
law. A successful production history would close one input to the dark
mass problem, not the gravitational and lensing completion.

\section{What changed and the next discriminating work}
The zero-temperature scalar bridge now has a sharper decision tree.
The original 1-MeV comparison cannot simply start from a thermally stable
empty dark origin and retain the earlier dilute cascade. Its intermediate
phase is a possible production route, but it requires its own non-equilibrium
field/particle calculation. A distinct sufficient-condition mass-family
route avoids that particular phase problem and can proceed to kinetic
production, while keeping its parameters explicitly selected.

Next execute the candidate's collision integrals with full charge channels,
finite-width intermediate states, thermal kinematics and number-changing
rates. For the symmetric mass-family branch, first verify thermal mass/rate
corrections and dark reprocessing before recalibrating abundance. For the
1-MeV phase branch, first derive the condensate and fluctuation evolution,
including the bath work and the changed spectrum. A stable equilibrium
table is not a substitute for either calculation.

Computational evidence: 108 bounded controls in the thermal/elastic run, with
five mass choices and nine temperatures, exact rational threshold identities,
stationary-root/log-grid comparisons, angular quadrature refinement, and
a time-dilated decay control. Its run 001 preceded the added decay control and
retains historical input hashes. The separate number-changing run 002
adds 120 exact/finite controls, including the independent topology sum;
its run 001 preceded that independent enumeration. The universal
sufficient-condition and amplitude derivations are written above; finite checks do not
prove physical validity. No simulations or external jobs remain running.

Literature scope: primary full text of \cite{seller} and the thermal section
of \cite{deromeri} inspected on 16 October 2026. A project-cache query for
the latter had no match. No PDFs retained. Broad singlet phase transitions,
thermal freeze-in and charge crossing are known; the present result is a
specific conditional screen for the current scalar construction. No complete
citation-graph search or global novelty claim is made.

\begin{thebibliography}{9}
\bibitem{seller} K. Seller, Z. Sz\'ep and Z. Tr\'ocs\'anyi,
\emph{Real effective potentials for phase transitions in models with extended
scalar sectors}, JHEP 04 (2023) 096, arXiv:2301.07961v2,
Appendix C.1, equations (C.11)--(C.12).
\url{https://arxiv.org/html/2301.07961v2}.
Only leading visible-bath thermal normalizations and the need to distinguish
them from a full resummed phase calculation are imported.
\bibitem{deromeri} V. De Romeri, D. Karamitros, O. Lebedev and T. Toma,
\emph{Neutrino dark matter and the Higgs portal: improved freeze-in analysis},
arXiv:2003.12606v2, sections 2.1.1 and 6.2.4 and Appendix A.
\url{https://arxiv.org/html/2003.12606v2}.
Its sterile-neutrino model is not identified with this scalar model;
the source supplies methodological overlap, not our abundance or phase verdict.
\end{thebibliography}
\end{document}
